%%%%%%%%%%%%%%%%%%%%%%%%
%
% $Autor: Wings $
% $Datum: 2020-07-24 09:05:07Z $
% $Pfad: GDV/Vortraege/latex - Ausarbeitung/Kapitel/Einleitung.tex $
% $Version: 4732 $
%
%%%%%%%%%%%%%%%%%%%%%%%%

\chapter{introduction}
An accurate Gait Pattern analysis can at the moment only be carriered out by health professionals  and is not available to individuales, even though that a healthy gait pattern reduces the risk of having back problems or even bad injuries especially for older poeple. 
analysing the gait pattern also gives us information about the person´s previous accidents or even about current neuronal problem as changes of gait patterns could in fact be a symptom of some sicknesses like Parkinson.
the long term goal of this project would be to anable individuales analyse their own gait pattern and recognise if there is something wrong with them at a reasonable price Tag. 
but for this paper the goal would be to use IMU sensors to collect information about the gait and identify if the person is walking, running or jumping.
As  hardware an Arduino Nano BLE 33  including an IMU Sensor will be used to collect and process the data, TinyML will be used because of the limited memory of the Embedded System, TensorFlow tools to process and load data (build or train the model.) TensorFlow Lite is then used to provide the ability to perform predictions on the trained model (load the model not to train the model). 
the tool created should be placed on the thigh in a pocket, as it was shows in AbhayasingheMurray:2014 anaysing one thigh is enough to accuratly detect the movement of one leg and to predict movements of the other one. IMU sensors should detect the angle of the thigh movement and it shoud then be saved on the hardware and analysed comparing the angle´s waveform to the previous ones registred during the research.
the hardware used has a limited memory, using tinyML helps us avoid this problem. 
Energy source is also a challenge, since the tool is to be carried arround, it has to have a portable source of energy. Using a battery could here be good choice since the device also needs to be small.
A power button needs to be added so that the tool doesn´t waste energy fruitlessly, as well as 4 LEDs to make the output understandable by the user: 
one for "Running" one for "Walking" one for "Jumping"  and one for "Sitting and Standing". this type of output could also be improved with using an LCD display instead of the 4 diffrent LEDs.
to make an accurate motion detection possible a database containing the diffrent 4 movements should be available, in this study the data used will be internal, taken by the same device that will then be used to detect the type of the motion and from diffrent poeple doing the same tasks for the same time period. having a good database makes data classification with machine learning better since specific patterns are repeated for each motion, particularities that a normal person couldn´t always point out. for that it is  important to achieve data homogeneity and to label all the informations recorded correctly. 
Along with a good database and to make an automatic classification of the gait pattern possible, deep learning will be needed to extract useful patterns from the data collected more precisely employing Neural networks (CNN), an algorithm inspired by the human brain that is commonly used to classify images, daily use of it is, for example, the face recognition feature on our smartphones, it is also be used for other needs such as speech recognition, fraud detection, medical diagnosis...
the library used in this project is Tensorflow an open-source software library for machine learning and artificial intelligence developed by the google brain Team. it is used to develop and train ML models but can not be used on mobile phones or embedded devices because of their limited memory, however, models can still be used on these devices with the help of TensorFlow Lite.

\section{Gangbildern}

\hfill \break
\textbf{das Gangbild}: bezeichnet die visuelle Abbildung der Bewegung und Haltung der Extremitäten sowie deren Zusammenspiel mit dem Rumpf während des Gangzyklus. 
Bei einem gesunden Gang bewegen sich die Extremitäten fließend und koordiniert, während die Haltung des Oberkörpers ausbalanciert wird. was normalerweise zu einer nach vorne gerichteten Verlagerung des Schwerpunkts führt. 

\hfill \break
\textbf{das Gangzyklus}: ist der Vorgang vom anfangs-Kontakt mit dem Boden bis zum nächsten anfangs-Kontak desselben Fusßes aus dem n\"achsten Zyklus. 
der Bodenkontakt zu Beginn des Gangzyklus ist der Anfang und wird als 0 \% des Gangzyklus bezeichnet. Der moment des nächsten kontakts desselben Fusßes wird als der 100\% punkt gennant .
Jeder Gangzykus wird zu zwei phasen erstmal unterteilt: Standphase und Schwungphase
Standphase: Periode in der der Fuß auf dem Boden ist (initial contact)
Schwungphase: Periode in der der Fuß in der Lunft ist und unter anderem der Schwung das Bein nach vorne bewegt diese phase beginnt mit dem Abheben des Fußes.

\hfill \break
 Ein bestimmte Gangbildmuster gibt es nicht denn jeder entwickelt seine eigene lösung , Menschen, die zum Beispiel in der Steppe Afrikas leben gehen mit anhaltender Knoegelenkflexion da die bodenverhältnisse diese besondere Anpassung der Lokomotion erfordern
Gangbild Unterschiede ergeben sich unter anderem aus folgenden Faktoren: Alter, Geschlecht,  K\"orpergr\"osße, K\"orperbau, Gewicht und Masseverteilung, Bodenbeschaffenheit, Lenemsumstände, Umgebzng... 

quelle %https://books.google.de/books?hl=de&lr=&id=fEqGIIg563UC&oi=fnd&pg=PA1&dq=gangzyklus+phasen+nach+perry&ots=QuprQttIG6&sig=6cB7tPIEenVKY3HD6KScpyVGNdU#v=onepage&q=gangzyklus%20phasen%20nach%20perry&f=false


 % ein bestimmte Gangbildmuster gibt es nicht denn jeder entwickelt seine eigene lösung : "Wie gelange ich mit einem Minimum an Anstrengung, mit ausreichend Stabilitat sowie einem guten Erscheinungsbild von einem Ort zum an- deren?"
\hfill \break
Abweichungen von einem normalen Gang, können Beschwerden verursachen und zu Verletzungen führen. 
 am meisten betroffen von Gangstörungen sind die \"U 60 da 10 \% der 60 j\"ahrigen und 30\% der 80 j\"ahrigen 
 Probleme beim gehen haben.  \hfill \break
 eine Gangbildanalyse hilft also abweichungen zu identifizieren um die nachcher zu behandeln und schwere Verlezungen vermeiden zu k\"onnen.  \hfill \break
die Identifizierung von abweichungen helfen auch Probleme bei der Patient zu erkenne und daher auch die Krankheit die es verursacht.  \hfill \break
ein Torkelnder Gang zum beispiel k\"onnte eine usache von Medikamente wie Schmerz- und Schlafmittel oder Antidepressiva die nicht für den Patienten geeignet sind, eine l\"osung daf\"ur w\"are mit dem behandelnden Arzt gekl\"art werden, ob das Medikament gewechselt werden kann

- Was uns eine gangbildanalyse bringt 
\section{TinyML}
\subsection{TinyML and Real Time}
\subsection{TinyMLWorkflow}
- Workflow 
- TensorFlow
- TensorFlow Lite 
- TensorFlow Micro 
\section{Arduino Nano 33 BLE}
\subsection{Hardware}
	\subsection{Software}
	- Datasheet lesen
	- Welcher Sensoren werden benutzt, warum und Welcher format von data ergeben die (IMU) 
	\section{Beschränkungen}
	- Beschränkungen identifizieren : Speicher, Gehäuse, Komplexität, Geschwindigkeit.   
	\section{IMU}
	\subsection{Einleitung}
	\subsection{Hardware}
	- Accelerometers
	- Gyroscope
	- Magnetometers
	\subsection{Software}
	\subsection{Calibration of Sensor Motion}
	\section{Arduino IDE}
	\subsection{Installation}
	\subsection{Configuration}
	\subsection{Tests}
	- Hello World - blink.ino
	-IMU
%\section{Aktuelles Market }
%- Medizinischer bereich
%- Everyday tools  
%\bigskip
